\documentclass[11pt]{article}
\usepackage[T1]{fontenc}
\usepackage[utf8]{inputenc}
\usepackage{lmodern,microtype,amsmath,amssymb,amsthm,mathtools}
\usepackage[margin=1in]{geometry}
\usepackage[colorlinks=true,linkcolor=blue,citecolor=blue,urlcolor=blue]{hyperref}
\setlength{\emergencystretch}{3em}
\newtheorem{theorem}{Theorem}[section]
\newtheorem{lemma}[theorem]{Lemma}
\newtheorem{proposition}[theorem]{Proposition}
\providecommand{\tightlist}{\setlength{\itemsep}{0pt}\setlength{\parskip}{0pt}}
\newcommand{\R}{\mathbb R}
\author{ChatGPT-6 Astra \and Kaarel H\"anni \and Dmitry Vaintrob\thanks{Authors are listed in alphabetical order.}}
\date{Research draft, 21 October 2026}
\title{Low-degree all-width learning bounds for shallow Bayesian neural networks}
\hypersetup{pdftitle={Low-degree all-width learning bounds for shallow Bayesian neural networks},pdfauthor={ChatGPT-6 Astra; Kaarel Hanni; Dmitry Vaintrob}}
\begin{document}
\maketitle
\begin{abstract}
For a tanh network with one hidden layer, we prove an all-width Gaussian
Bayesian learning bound essentially linear in the weighted squared norm
of a small teacher. The result includes the raw Gaussian prior, the prior
conditioned on an absolute margin everywhere, and a specified robust prior.
Training requires signed outputs above a fixed positive margin. The sample
bound is independent of learner width and holds at every mean-field or
stronger Gaussian scale. With hidden coefficients at most the teacher width
and bounded readout coefficients, the width dependence is quadratic.
Allowing the readout coefficients also to reach the width gives a quartic
sufficient bound. The proof represents tanh ridges as
mixtures of halfspaces, compresses their votes, and controls posterior
readout energy. Independent AI-agent review checked the argument; it is
formalized in Lean for the main learning statements, with scope recorded in the repository.
\end{abstract}
\section{Result and scope}\label{result-and-scope}

For one hidden layer, the mean-field-and-stronger Gaussian learner
admits a direct finite-class argument. Its sample bound is linear, up to
logarithms, in a normalized target-cylinder cost. No deep-network
compression is needed.

Let \(d,C\ge1\), put \(\bar x=(x,1)\), and write the learner as
\(z_w(x)=\sum_{r=1}^N v_r\tanh(W_r\cdot\bar x)\). For a teacher

\[
z_*(x)=\sum_{j=1}^{C}a_j\tanh(b_j\cdot\bar x),\qquad
f(x)z_*(x)>\tfrac12\quad(x\in\{-1,1\}^d),
\]

write \(A_1=\sum_j|a_j|\), \(A_2^2=\sum_j a_j^2\), and
\(U=d\min\{\max_j\|b_j\|_2^2,(2/C)\sum_j\|b_j\|_2^2\}\). The learner has
arbitrary width \(N\ge C\), independent Gaussian entries of variance
\(t^2/d\) in the first layer and \(t^2/N^2\) in the readout, where
\(0<t\le1\).

Any of the following priors is allowed:

\begin{enumerate}
\def\labelenumi{\arabic{enumi}.}
\tightlist
\item
  The Gaussian prior conditioned on \(|z(x)|>1/2\) at every Boolean
  input.
\item
  A robustness-restricted prior, using one hidden layer, radius
  \(c/\sqrt{1+E_v}\), \(E_v=N\|v\|_2^2\), clipped-score tolerance
  \(1/16\), and optional additional global-margin conditioning.
\item
  The raw Gaussian prior, with no initial margin or robustness
  restriction.
\end{enumerate}

Here the shallow robustness event means the following. Put
\(T(u)=\max\{-1/2,\min\{1/2,u\}\}\), add disturbances by
\(z_w(x;\eta)=v^\top[\tanh(W\bar x)+\eta^1]+\eta^{\rm out}\), and let
\(\|\eta\|_*^2=\|\eta^1\|_2^2/N+|\eta^{\rm out}|^2\).
The event requires
\[
 |T(z_w(x;\eta))-T(z_w(x))|\le1/16
 \quad\text{for every }x\in\{-1,1\}^d,
 \quad\|\eta\|_*\le c/\sqrt{1+N\|v\|_2^2}.
\]
The calibration \(c>0\) is fixed before training.

Fix an arbitrary input distribution \(D\) on the Boolean cube and draw the inputs independently from \(D\). In all three cases, training conditions on \(f(X_i)z(X_i)>1/2\). Test
error means \(f(X)z(X)\le1/4\), which bounds ordinary sign error. With
global-margin conditioning, the two errors coincide; this equality is
not asserted for the raw or robust-only priors.

Let \(s=3\) in cases 1 and 3, and \(s=2c+3\) in case 2. Set

\[
\rho=\frac1{4[1+sA_1(d+1)/\sqrt d]},\qquad
D_*=U+2s^2C A_2^2+(d+2)\log\frac{e\sqrt{2\pi}}{2\rho}.
\]

Then \(D_*>1\). For \(0<\varepsilon,\delta<1\), define

\[
B^2=8\left[D_*+\log\frac4\varepsilon\right],\qquad
m=\left\lceil512 B^2\log\frac{24}{\varepsilon}\right\rceil,
\qquad h_d=(d+1)(1+d\log2).
\]

A sufficient sample size is

\[
\boxed{\quad n\ge\frac{16}{\varepsilon}
\left[m h_d+\log\frac1\delta\right].\quad}
\]

With probability at least \(1-\delta\), the posterior-averaged
test-margin error is at most \(\varepsilon\), simultaneously for all
\(N\ge C\) and \(0<t\le1\).

Thus the bound is

\[
O\!\left(\frac{d^2[D_*+\log(1/\varepsilon)]\log(1/\varepsilon)
 +\log(1/\delta)}{\varepsilon}\right).
\]

A sharper alternative, proved below by a relative growth-function
argument, is obtained by setting

\[
T=\frac{64}{\varepsilon}
\left[m(d+1)+\log\frac{4(m+2)}{\delta}\right]
\]

and taking \(n\ge\lceil4T\log(8eT)\rceil\). Up to logarithms, this is
\(\widetilde O((dD_*+\log(1/\delta))/\varepsilon)\), improving the
dimension factor from \(d^2\) to \(d\).

The displayed constants are deliberately loose; the useful improvement
is in the powers.

In terms of the teacher's weighted squared norm

\[
\kappa_C=\frac dC\sum_j\|b_j\|_2^2+C\sum_j a_j^2,
\]

one has \(D_*\le2s^2\kappa_C+(d+2)\log(e\sqrt{2\pi}/(2\rho))\). Thus the
sample bound is essentially \textbf{linear in the teacher's weighted
squared norm}.

\subsection{Width consequences}\label{width-consequences}

If first-layer entries satisfy \(|b_{ji}|\le W\) and readout entries
satisfy \(|a_j|\le V\), then

\[
D_*\le d(d+1)W^2+2s^2 C^2V^2
 +(d+2)\log\!\left(2e\sqrt{2\pi}[1+sCV(d+1)/\sqrt d]\right).
\]

\begin{itemize}
\tightlist
\item
  For fixed \(V\), with \(W\le C\), this gives a genuinely quadratic
  dependence on teacher width \(C\). The first, finite-cube bound is
  literally \(O(C^2)\) when the other parameters are fixed: its additive
  \(\log C\) target-cost term is absorbed into \(C^2\). The
  dimension-sharpened bound has an additional logarithmic factor in
  \(C\). Both statements hold even for the raw Gaussian prior, without
  an initial margin or robustness restriction.
\item
  Under the original bounded-weight-complexity convention \(W,V\le C\),
  it gives \(D_*=O(d^2C^2+s^2C^4+d\log(dCs))\), and hence quartic
  dependence on \(C\). This proof does \textbf{not} establish quadratic
  dependence for that full class.
\item
  Width alone cannot be substituted for the theorem's teacher resources
  without controlling coefficient magnitudes or another implementation
  cost.
\end{itemize}

\section{Proof}\label{proof}

\subsection{1. A target cylinder at every
width}\label{a-target-cylinder-at-every-width}

Partition the \(N\) hidden neurons into \(C\) balanced groups with
positive sizes \(k_j\). Then \(N/k_j\le2C\) and \(k_j/N\le2/C\). Give
each member of group \(j\) first-layer row \(b_j\) and readout
coefficient \(sa_j/k_j\). This clone has output \(sz_*\).

The Gaussian coordinates are \(\sqrt d\,W_{ri}\) and \(Nv_r\), all
independent \(N(0,t^2)\). Require every coordinate to lie within
\(\rho t\) of its clone center. On this cylinder,

\[
|z_w(x)-sz_*(x)|
\le \rho+sA_1\rho(d+1)/\sqrt d=\tfrac14
\]

simultaneously on the entire Boolean cube. The first term is the readout
perturbation, using bounded activations; the second is the changed
hidden activations weighted by the cloned readout.

In cases 1 and 3 this gives \(f(x)z_w(x)>5/4\). In case 2 it gives
\(f(x)z_w(x)>c+5/4\).

For a shallow network, a simultaneous hidden-and-output disturbance
obeys

\[
|z_w(x;\eta)-z_w(x)|
=|v\cdot\eta^1+\eta^{\rm out}|
\le\sqrt{1+E_v}\,\|\eta\|_*.
\]

Consequently the calibrated ball changes any output by at most \(c\).
Every cylinder network in case 2 remains above signed margin \(5/4\)
throughout that ball. Its clipped score is constant there, so it lies in
the robust support. (In fact, for \(c\le1/16\), every shallow network
lies in that support, because clipping is 1-Lipschitz.)

There are \(N(d+2)\) independent constrained coordinates. Their squared
centers sum to at most

\[
N U+s^2N^2\sum_j a_j^2/k_j\le N[U+2s^2C A_2^2].
\]

For \(Z\sim N(0,t^2)\), \(0<\rho\le1/4\),

\[
\Pr(|Z-u|\le\rho t)
\ge\frac{2\rho}{e\sqrt{2\pi}}e^{-u^2/t^2}.
\]

Multiplying, and using \(t\le1\), proves that the cylinder has raw mass
at least \(e^{-ND_*/t^2}\). It lies in every correctly labeled
posterior-conditioning event. In case 3 that event consists only of the
training-margin constraints; no initial support restriction is used.

\subsection{2. A width-independent readout-energy
tail}\label{a-width-independent-readout-energy-tail}

Under the raw Gaussian, \(NE_v/t^2\) has the chi-square distribution
with \(N\) degrees of freedom. For \(R\ge8\), its exponential moment
gives

\[
\Pr(E_v>R)\le e^{-NR/(8t^2)}.
\]

Dividing by the posterior denominator from step 1, and taking \(R=B^2\),
shows that the posterior mass of \(E_v>B^2\) is at most
\(\varepsilon/4\), on every sample. Every retained network has

\[
\sum_{r=1}^N|v_r|\le\sqrt N\,\|v\|_2\le B.
\]

Let \(Q'\) be the posterior conditioned on this readout-energy event.

\subsection{3. A bounded-variation shallow network is a mixture of
halfspaces}\label{a-bounded-variation-shallow-network-is-a-mixture-of-halfspaces}

Let \(T\) have continuous distribution function
\(\Pr(T\le u)=(1+\tanh u)/2\). Then

\[
\tanh(u)=\mathbb E_T\operatorname{sign}(u-T).
\]

Thus each ridge \(\tanh(b\cdot\bar x)\) is an average of Boolean affine
halfspaces on the cube. A sign change of a halfspace is another
halfspace pattern; zero-probability boundary ties can be ignored because
the input set is finite.

Every retained output \(z\) therefore satisfies

\[
z(x)=B\,\mathbb E H(x)
\]

for a distribution on Boolean affine-halfspace patterns \(H\). Assign
mass \(|v_r|/B\) to the corresponding signed ridge mixture. Fill any
remaining mass with equal probabilities on the two constant halfspaces.

The number \(M_d\) of distinct affine-halfspace patterns on the cube is
at most

\[
M_d\le(e2^d)^{d+1},\qquad\log M_d\le h_d,
\]

by the VC dimension \(d+1\) and Sauer's lemma.

Independently sample \(m\) halfspaces from the representation of \(z\)
and put

\[
h(x)=\frac Bm\sum_{j=1}^m H_j(x).
\]

For every fixed input \(x\), Hoeffding's inequality gives

\[
\Pr(|h(x)-z(x)|>1/16)
\le2e^{-m/(512B^2)}\le\alpha:=\varepsilon/12.
\]

The compressed outputs lie in one fixed finite family \(\mathcal H\),
with \(\log|\mathcal H|\le m h_d\). Its sampling distribution may depend
on \(z\), but does not depend on the input being tested. Averaging these
distributions over \(w\sim Q'\) gives a common law \(\Gamma\) on
\(\mathcal H\).

\subsection{4. Generalization of the compressed
mixture}\label{generalization-of-the-compressed-mixture}

Use the ramp loss

\[
\psi(u)=\begin{cases}
1,&u\le5/16,\\
(7-16u)/2,&5/16<u<7/16,\\
0,&u\ge7/16.
\end{cases}
\]

On a training point, a successful approximation has signed compressed
output strictly above \(7/16\), because every original posterior draw
has signed output above \(1/2\). Therefore

\[
\mathbb E_{h\sim\Gamma}\widehat R_\psi(h)\le\alpha.
\]

At a test input where \(fz_w\le1/4\), successful approximation gives
\(fh\le5/16\), hence unit ramp loss. Thus

\[
R_{1/4}(Q')\le\mathbb E_{h\sim\Gamma}R_\psi(h)+\alpha.
\]

For a finite family of bounded losses, the relative Chernoff inequality
and a union bound imply, with probability at least \(1-\delta\),

\[
R_\psi(h)\le2\widehat R_\psi(h)
+\frac{8[\log|\mathcal H|+\log(1/\delta)]}{n}
\quad\text{for every }h\in\mathcal H.
\]

This event is independent of \(N,t\), and holds for all mixture laws on
\(\mathcal H\), including the data-dependent law \(\Gamma\). Hence

\[
R_{1/4}(Q)\le\varepsilon/4+3\alpha
+\frac{8[mh_d+\log(1/\delta)]}{n}\le\varepsilon.
\]

This proves the simultaneous theorem.

\section{Checks and limitations}\label{checks-and-limitations}

\begin{itemize}
\tightlist
\item
  The proof controls the posterior average, so tied minimizers and
  posterior mixtures cause no problem. A correct-margin posterior
  confidence threshold \(q<1\) follows by using risk tolerance
  \((1-q)\varepsilon\).
\item
  The arbitrary first-layer weights of posterior draws do not enter the
  compressed family. Only the halfspace patterns matter.
\item
  The raw-prior version still uses hard training-margin conditioning. It
  does not assert the same result for Gaussian regression, logistic
  likelihoods, or conditioning only on training signs without a positive
  margin.
\item
  The theorem works with the existing Gaussian mean-field and stronger
  scalings: \(t=N^{(1-\gamma)/2}\) gives every \(\gamma\ge1\). It also
  covers every sufficiently strong finite-width scale of the original
  strong-prior parametrization, by setting \(t=\sqrt N\sigma\le1\).
\item
  The first bound uses simple finite-cube counting. The additional
  argument below improves the dimension factor to \(d\), up to
  logarithms, by counting the sample patterns of voted halfspaces.
\item
  A Boolean circuit with a single output gate and only one Boolean stage
  is a small target class regardless of the number of unused wires. The
  meaningful shallow comparator here is a real-valued tanh hidden layer
  and linear readout implementing a Boolean function.
\item
  This is an independent derivation, not a claimed optimization of the
  old degree-510/10974 compression formulas.
\end{itemize}

\section{Sharpening the dimension factor by counting sample
patterns}\label{sharpening-the-dimension-factor-by-counting-sample-patterns}

The preceding proof counted halfspaces on the entire cube. Counting only
their patterns on a sample improves the dimension factor from \(d^2\) to
\(d\), up to logarithms. Here is a self-contained argument with
deliberately generous constants.

\subsection{A relative growth-function
lemma}\label{a-relative-growth-function-lemma}

Let \(\mathcal A\) be a finite class of measurable sets and let
\(\Pi_{\mathcal A}(r)\) be its maximum number of distinct membership
patterns on \(r\) points. For an iid sample of size \(n\), with
probability at least \(1-\delta\),

\[
P(A)\le2P_n(A)+\frac{32}{n}
\log\frac{4\Pi_{\mathcal A}(2n)}{\delta}
\quad\text{for all }A\in\mathcal A.
\]

To prove this, let \(u>0\) with \(nu\ge32\log2\), and suppose some set
has \(p=P(A)>2P_n(A)+u\). Conditional on the first sample and a selected
witness \(A\), an independent ghost sample satisfies \(P'_n(A)\ge3p/4\)
with probability at least \(1/2\): the lower Chernoff tail is at most
\(e^{-np/32}\). Hence the probability of a witness is at most twice the
probability that

\[
P'_n(A)>\tfrac32P_n(A)+\tfrac34u
\]

for some \(A\).

Pool the two samples and split their \(2n\) positions uniformly into two
groups of size \(n\). A fixed membership pattern with \(k\) positive
positions has a hypergeometric number \(K\) of positives in the first
group, with mean \(k/2\). The displayed event implies

\[
K<\tfrac25k-\tfrac3{10}nu,
\]

so the sampling-without-replacement form of Hoeffding's inequality
gives, for \(k>0\),

\[
\Pr\left(K<\tfrac25k-\tfrac3{10}nu\right)
\le \exp\!\left[-\frac{2(k/10+3nu/10)^2}{k}\right]
\le e^{-6nu/25}.
\]

For the last inequality, expand the exponent and use
\(k/50+9(nu)^2/(50k)\ge3nu/25\). The case \(k=0\) is impossible. A union
bound over at most \(\Pi_{\mathcal A}(2n)\) patterns therefore bounds
the witness probability by \(2\Pi_{\mathcal A}(2n)e^{-6nu/25}\). Taking
\(u=(32/n)\log[4\Pi_{\mathcal A}(2n)/\delta]\) proves the lemma.

\subsection{Applying the lemma to the compressed ramp
losses}\label{applying-the-lemma-to-the-compressed-ramp-losses}

Work on labeled examples \((x,y)\). Consider all sets

\[
\{(x,y):\psi(yh(x))>q\},\qquad h\in\mathcal H,\quad 0\le q\le1.
\]

This is a finite set class: each compressed output is a vote of \(m\)
Boolean halfspaces and has at most \(m+1\) possible output values, so
varying \(q\) gives at most \(m+2\) membership patterns for a fixed
vote.

On any \(r\) input points, one affine halfspace has at most
\((er)^{d+1}\) patterns. Choosing \(m\) such patterns gives at most
\((er)^{m(d+1)}\) voted outputs. Multiplication by the fixed labels of
the sample does not increase this number. Consequently the
loss-level-set growth function satisfies

\[
\Pi(r)\le(m+2)(er)^{m(d+1)}.
\]

Apply the relative growth-function lemma to all these level sets and
integrate in \(q\). The layer-cake identity for a bounded nonnegative
loss gives, simultaneously for every \(h\in\mathcal H\),

\[
R_\psi(h)\le2\widehat R_\psi(h)
+\frac{32}{n}\left[m(d+1)\log(2en)
+\log\frac{4(m+2)}{\delta}\right].
\]

The same averaging and readout-tail argument as above therefore gives
posterior risk at most \(\varepsilon\) whenever

\[
\boxed{\quad n\ge\frac{64}{\varepsilon}
\left[m(d+1)\log(2en)+\log\frac{4(m+2)}{\delta}\right].\quad}
\]

For a fully explicit sufficient bound, put

\[
T=\frac{64}{\varepsilon}
\left[m(d+1)+\log\frac{4(m+2)}{\delta}\right].
\]

Then it suffices to take

\[
\boxed{\qquad n\ge\left\lceil4T\log(8eT)\right\rceil.\qquad}
\]

Indeed, write \(A=64m(d+1)/\varepsilon\) and
\(B_0=64\log[4(m+2)/\delta]/\varepsilon\), so \(T=A+B_0\). At
\(n_0=4T\log(8eT)\), the quantity \(A\log(2en_0)+B_0\) is at most
\(T[\log(8eT)+\log\log(8eT)+1]\le n_0\). The difference
\(n-A\log(2en)-B_0\) is increasing for \(n\ge n_0>A\).

Up to logarithmic factors, this yields the sharper sample bound

\[
\boxed{\quad n=\widetilde O\!\left(
\frac{dD_*+\log(1/\delta)}{\varepsilon}\right).\quad}
\]

The exact formulas above include the logarithms in \(\varepsilon\),
\(D_*\), \(d\), and confidence. The sample event remains independent of
learner width and Gaussian scale.

\subsection{References for the elementary
ingredients}\label{references-for-the-elementary-ingredients}

\begin{itemize}
\tightlist
\item
  W. Hoeffding, ``Probability inequalities for sums of bounded random
  variables,'' \emph{Journal of the American Statistical Association}
  \textbf{58} (1963), 13--30. The sampling-without-replacement
  inequality is included in Section 6.
\item
  N. Sauer, ``On the density of families of sets,'' \emph{Journal of
  Combinatorial Theory, Series A} \textbf{13} (1972), 145--147.
\end{itemize}

The relative growth-function statement used here was proved above,
rather than being assumed as an external theorem.

\paragraph{Acknowledgments and status.}
The parameter conventions and prior-conditioning questions come from the
collaborative NNBayes research program. This paper supplies a separate
shallow proof; the accompanying review records describe its independent
AI-agent mathematical checks. The paper has not received a full human
mathematical review.
\end{document}
